% Part: lambda-calculus
% Chapter: syntax
% Section: free-variables

\documentclass[../../../include/open-logic-section]{subfiles}

\begin{document}

\olfileid{lam}{syn}{fv}
\olsection{Variabel Bebas}

Kalkulus lambda ngrembug fungsi, lan
abstraksi lambda iku cara mbentuk fungsi.
Miturut intuisi, $\lambd[x][M]$ iku
fungsi kang bijine diwenehake dening~$M$
nalika argumen fungsi iku diwenehake
marang~$x$. Nanging ora saben muncule~$x$
ing~$M$ nduweni peran kanggo abstraksi iki:
yen $M$ ngemot abstraksi liyane
$\lambd[x][N]$, mula muncule $x$ ing~$N$
duweni peran kanggo $\lambd[x][N]$
nanging ora kanggo $\lambd[x][M]$.
Dadi, abstraksi lambda $\lambd[x]$ ing
$\lambd[x][M]$ \emph{ngiket} muncule
$x$ ing~$M$ kang durung diiket dening
abstraksi lambda liyane, yaiku muncule
$x$ kang \emph{bebas} ing~$M$.

\begin{defn}[Cakupan]
Yen $\lambd[x][M]$ muncul ing suku~$N$,
mula muncule~$M$ kang cocog iku
\emph{cakupan} saka~$\lambd[x]$.
\end{defn}


\begin{defn}[Muncule variabel bebas lan kaiket]
  Muncule variabel $x$ ing suku~$M$
  diarani \emph{bebas} yen ora ana
  ing cakupan sawijining $\lambd[x]$,
  lan \emph{kaiket} yen ana.
  Muncule variabel~$x$ ing
  $\lambd[x][M]$ diiket dening
  $\lambd[x]$ ing wiwitan yen lan
  mung yen muncule $x$ ing~$M$
  bebas.
\end{defn}

\begin{ex}
  Ing $\lambd[x][x y]$, $x$ lan $y$
  loro-lorone ana ing cakupan
  $\lambd[x]$, mula $x$ diiket
  dening $\lambd[x]$. Amarga $y$
  ora ana ing cakupan sembarang
  $\lambd[y]$, variabel iku bebas.
  Ing $\lambd[x][x x]$, loro
  muncule $x$ diiket dening~$\lambd[x]$,
  amarga loro-lorone bebas ing $xx$.
  Ing $((\lambd[x][xx])x)$,
  muncule~$x$ kang pungkasan bebas,
  amarga ora ana ing cakupan
  $\lambd[x]$. Ing
  $\lambd[x][(\lambd[x][x])x]$,
  cakupan $\lambd[x]$ kapisan yaiku
  $(\lambd[x][x])x$ lan cakupan
  $\lambd[x]$ kapindho yaiku muncule~$x$
  kang sadurunge pungkasan.
  Ing $(\lambd[x][x])x$, muncule
  $x$ kang pungkasan bebas, dene
  muncule sadurunge pungkasan kaiket.
  Dadi, muncule $x$ kang sadurunge
  pungkasan ing
  $\lambd[x][(\lambd[x][x])x]$
  diiket dening $\lambd[x]$ kapindho,
  dene muncule kang pungkasan
  diiket dening $\lambd[x]$ kapisan.
\end{ex}

Kanggo suku $P$, kita bisa mriksa
kabeh muncule variabel ing njero
lan nemtokake himpunan variabel
bebas. Himpunan iki dilambangake
$\FV{P}$ lan nduweni definisi
lumrah kaya mangkene:

\begin{defn}[Variabel bebas ing sawijining suku] \ollabel{def:fv}
  Himpunan \emph{variabel bebas}
  saka suku ditetepake kanthi
  induktif kaya mangkene:
  \begin{enumerate}
    \item $\FV{x} = \{x\}$ \ollabel{def:fv1}
    \item $\FV{\lambd[x][N]} = \FV{N} \setminus \{x\}$ \ollabel{def:fv2}
    \item $\FV{PQ} = \FV{P} \cup \FV{Q}$ \ollabel{def:fv3}
  \end{enumerate}
\end{defn}

\begin{prob}
  \begin{enumerate}
  \item Temtokna cakupan $\lambd[g]$
    lan loro $\lambd[x]$ ing suku
    iki: $\lambd[g][(\lambd[x][g (x x)]) \lambd[x][g (x x)]]$.
  \item Ing $\lambd[g][(\lambd[x][g (x x)]) \lambd[x][g (x x)]]$,
    apa kabeh muncule variabel kaiket?
    Abstraksi endi kang ngiket
    saben muncule iku?
  \item Temtokna
    $\FV{\lambd[x][(\lambd[y][(\lambd[z][x y]) z]) y]}$.
  \end{enumerate}
\end{prob}

\begin{explain}
Variabel bebas iku kaya rujukan
marang donya njaba, yaiku
\emph{lingkungan}; suku kang
ngemot variabel bebas bisa
dianggep minangka suku kang
durung ditemtokake kanthi
lengkap, amarga tumindake
gumantung marang carane kita
netepake lingkungan. Umpamane,
ing suku $\lambd[x][f x]$,
kang nampa argumen~$x$ lan
ngasilake asil $f$ ing argumen
iku, variabel~$f$ bebas.
Biji suku iki gumantung marang
lingkungan, mligine biji $f$
ing lingkungan kasebut.

Yen suku iki diabstraksi maneh,
kita oleh $\lambd[f][\lambd[x][f
    x]]$. Suku iki ora gumantung
maneh marang variabel lingkungan
$f$, amarga saiki makili fungsi
kang nampa rong argumen lan
ngasilake asil aplikasi argumen
kapisan marang argumen kapindho.
Ngganti $f$ ing lingkungan
ora bakal ngowahi tumindake
suku iki, amarga suku mung
nggunakake apa kang diwenehake
minangka argumen, dudu biji
$f$ ing lingkungan.
\end{explain}

\begin{defn}[Suku katutup, kombinator]
  Suku kang ora nduweni
  variabel bebas diarani
  \emph{suku katutup}, utawa
  \emph{kombinator}.
\end{defn}

\begin{lem}\ollabel{lem:fv}
  \begin{enumerate}
    \item \ollabel{lem:fv-abs}
      Yen $y \neq x$, mula
      $y \in
      \FV{\lambd[x][N]}$ yen lan mung
      yen $y \in \FV{N}$.
    \item \ollabel{lem:fv-app}
      $y \in \FV{PQ}$ yen lan mung yen
      $y \in \FV{P}$ utawa
      $y \in \FV{Q}$.
  \end{enumerate}
\end{lem}

\begin{proof}
  Latihan.
\end{proof}

\begin{prob}
  Buktikna \olref[lam][syn][fv]{lem:fv}.
\end{prob}

\end{document}
